\documentclass[UTF8,a4paper,11pt]{ctexart}

% ── Page geometry ──
\usepackage[margin=2.2cm,top=2.5cm,bottom=2.5cm]{geometry}

% ── Typography ──
\usepackage{fontspec}
\usepackage{setspace}
\setstretch{1.2}

% ── Graphics ──
\usepackage{graphicx}
\graphicspath{{/home/zhiwen/projects/mooc2handout-skill/coursera-mcp/module-2/frames/}}

% ── Colors ──
\usepackage{xcolor}
\definecolor{accent}{HTML}{1a73e8}
\definecolor{darkgray}{HTML}{333333}
\definecolor{lightgray}{HTML}{f5f5f5}
\definecolor{keyframebg}{HTML}{fafafa}
\definecolor{videoblue}{HTML}{065fd4}
\definecolor{notebg}{HTML}{e8f0fe}
\definecolor{warnbg}{HTML}{fef7e0}

% ── Links ──
\usepackage{hyperref}
\hypersetup{
  colorlinks=true,
  linkcolor=accent,
  urlcolor=videoblue,
  citecolor=accent,
  pdftitle={模型上下文协议导论},
  pdfauthor={}
}

% ── Layout ──
\usepackage{enumitem}
\setlist[itemize]{leftmargin=1.8em,itemsep=2pt}
\setlist[enumerate]{leftmargin=1.8em,itemsep=2pt}
\usepackage{booktabs}
\usepackage{longtable}
\usepackage{tcolorbox}
\usepackage{titlesec}
\usepackage{fancyhdr}
\usepackage{lastpage}
\usepackage{amssymb}

% ── Header/Footer ──
\pagestyle{fancy}
\fancyhf{}
\fancyhead[L]{\small\textcolor{darkgray}{模型上下文协议导论}}
\fancyhead[R]{\small\textcolor{darkgray}{第二单元：搭建服务器与工具}}
\fancyfoot[C]{\small\textcolor{darkgray}{\thepage\ / \pageref{LastPage}}}
\renewcommand{\headrulewidth}{0.4pt}

% ── Section styling ──
\titleformat{\section}
  {\Large\bfseries\color{accent}}
  {\thesection}{1em}{}
\titleformat{\subsection}
  {\large\bfseries\color{darkgray}}
  {\thesubsection}{1em}{}

% ── Custom environments ──
\newtcolorbox{keyconcept}{
  colback=notebg, colframe=accent,
  fonttitle=\bfseries, boxrule=0.5pt,
  left=8pt, right=8pt, top=6pt, bottom=6pt
}
\newtcolorbox{supplement}{
  colback=lightgray, colframe=darkgray,
  fonttitle=\bfseries, boxrule=0.3pt,
  left=8pt, right=8pt, top=6pt, bottom=6pt
}
\newtcolorbox{exercise}{
  colback=warnbg, colframe=orange,
  fonttitle=\bfseries, boxrule=0.5pt,
  left=8pt, right=8pt, top=6pt, bottom=6pt
}

% ── Video link command ──
\newcommand{\videolink}[2]{%
  \href{#1}{\textcolor{videoblue}{\textbf{#2}}}%
}

% ── Keyframe figure command ──
\newcommand{\keyframe}[2]{%
  \begin{center}
    \fcolorbox{lightgray}{keyframebg}{%
      \includegraphics[width=0.75\linewidth]{#1}%
    }
    \par\vspace{2pt}
    {\small\textcolor{darkgray}{#2}}
  \end{center}
}

\begin{document}

\begin{center}
  \vspace*{2cm}
  {\Huge\bfseries\color{accent} 模型上下文协议导论}\\[0.3cm]
  {\small\textcolor{darkgray}{Introduction to Model Context Protocol}}\\[0.8cm]
  {\LARGE 第二单元：搭建服务器与工具}\\[0.3cm]
  {\small\textcolor{darkgray}{Module 2: Setting Up the Server and Tools}}\\[1.2cm]
  {\large 授课教师： }\\[0.5cm]
  {\large\textcolor{darkgray}{由 mooc2handout 自动生成}}\\[0.3cm]
  {\small\textcolor{darkgray}{\today}}
  \vfill
  {\small 本讲义由课程字幕、补充材料与可选的 AI 推断关键帧自动生成。}
\end{center}
\newpage

\tableofcontents
\newpage

\section*{本单元知识地图}
\addcontentsline{toc}{section}{本单元知识地图}
\begin{longtable}{p{3.1cm}p{10.2cm}}
\toprule
主题 & 核心问题 \\
\midrule
练习项目 & 为什么要先做一个 CLI 聊天机器人？ \\
工具定义 & read / update 文档分别对应什么操作？ \\
调试闭环 & Inspector 在开发流程里扮演什么角色？ \\
\bottomrule
\end{longtable}

\section*{0. 概要}
\addcontentsline{toc}{section}{0. 概要}
这一单元进入动手阶段，围绕一个 CLI 聊天机器人项目搭建 MCP server 并验证工具。
核心目标不是写出炫技代码，而是把“定义能力 -> 运行调试 -> 修正实现”的闭环建立起来。

\begin{itemize}
  \item 先用一个小项目把 MCP 落地。
  \item 工具是 server 暴露给模型的可调用动作。
  \item Inspector 能让你快速检查 server 的实际行为。
\end{itemize}

\section*{1. 项目搭建}
\addcontentsline{toc}{section}{1. 项目搭建}
\noindent\textbf{回看：}\videolink{https://www.coursera.org/learn/introduction-to-model-context-protocol/home/module/2}{项目搭建}

\textbf{本讲重点：} 先把练习项目搭起来，再往里填 MCP 逻辑。

这一讲启动了一个 CLI 聊天机器人项目，用它来观察 client 和 server 如何真正协作。
项目里会用到一组存放在内存中的文档，方便我们把注意力集中在 MCP 的结构，而不是持久化细节。

这个设计的好处是：你可以一边在终端里聊天，一边逐步把服务器能力补进去，
每一个新功能都能马上看到效果。

\begin{itemize}
  \item 这是一个练习性质的 CLI 聊天机器人，而不是完整产品。
  \item 文档先放在内存里，便于快速迭代。
  \item 项目目标是让你看到 client 和 server 如何协同工作。
\end{itemize}

\begin{keyconcept}
\textbf{自动摘要}
根据字幕自动扩写，这一讲再补充几条最值得记住的线索。
\begin{itemize}
  \item \textbf{形式定义}: 这一讲梳理 AND / OR / NOT / IF...THEN 等连接词的真值条件，并提醒不要把日常语言的直觉直接搬进数学里。
\end{itemize}
\end{keyconcept}

\section*{2. 使用 MCP 定义工具}
\addcontentsline{toc}{section}{2. 使用 MCP 定义工具}
\noindent\textbf{回看：}\videolink{https://www.coursera.org/learn/introduction-to-model-context-protocol/home/module/2}{使用 MCP 定义工具}

\textbf{本讲重点：} 开始给 server 增加两个最基础的 tool。

讲者在 \texttt{mcpserver.py} 里准备了一个基础 MCP server，
接下来要补上的就是两个工具：读取文档和更新文档。
这一步的重点不是语法炫技，而是理解 tool 是 server 暴露出来、供模型调用的能力。

这个单元还提醒你，MCP tool 的定义通常要带上清晰的输入输出结构；
从工程角度看，这就是把“能做什么”和“怎么调用”拆成可验证的接口。

\begin{itemize}
  \item 先实现 read / update 这样的基础能力。
  \item 工具代码放在 	exttt{mcpserver.py} 中。
  \item JSON schema 让工具接口更明确、可检查。
\end{itemize}

\begin{keyconcept}
\textbf{自动摘要}
根据字幕自动扩写，这一讲再补充几条最值得记住的线索。
\begin{itemize}
  \item \textbf{形式定义}: 这一讲梳理 AND / OR / NOT / IF...THEN 等连接词的真值条件，并提醒不要把日常语言的直觉直接搬进数学里。
\end{itemize}
\end{keyconcept}

\section*{3. 服务器 Inspector}
\addcontentsline{toc}{section}{3. 服务器 Inspector}
\noindent\textbf{回看：}\videolink{https://www.coursera.org/learn/introduction-to-model-context-protocol/home/module/2}{服务器 Inspector}

\textbf{本讲重点：} 用 Inspector 把 server 的行为看清楚。

这一讲介绍如何运行 \texttt{mcp dev}，再在浏览器里打开 MCP inspector。
这套工具把调试回路缩短了：你能直接连接 server，切换资源、提示词和工具标签页，
观察它到底返回了什么。

讲者也提醒，Inspector 仍在快速演进，所以界面细节可能变化。
但核心思路不会变：先连上 server，再逐项验证能力是否符合预期。

\begin{itemize}
  \item Inspector 是验证 server 是否按预期工作的调试面板。
  \item 它能帮助你尽早发现工具定义或返回值的问题。
  \item 把开发、测试和修正连成一个短闭环很重要。
\end{itemize}

\begin{keyconcept}
\textbf{自动摘要}
根据字幕自动扩写，这一讲再补充几条最值得记住的线索。
\begin{itemize}
  \item \textbf{形式定义}: 这一讲梳理 AND / OR / NOT / IF...THEN 等连接词的真值条件，并提醒不要把日常语言的直觉直接搬进数学里。
\end{itemize}
\end{keyconcept}

\section*{4. 本单元最该记住的五句话}
\addcontentsline{toc}{section}{4. 本单元最该记住的五句话}
\begin{enumerate}
  \item 项目先行，概念才能落地。
  \item tools 是 server 端最直接的能力暴露方式。
  \item mcp dev + Inspector 构成了很实用的调试闭环。
  \item 把文档放进内存可以暂时降低复杂度。
  \item 清晰的接口比复杂的实现更重要。
\end{enumerate}

\section*{5. 练习题}
\addcontentsline{toc}{section}{5. 练习题}
\subsection*{A. 选择题（5题）}
\begin{enumerate}
  \item 本单元要做的项目是什么？\\
  A. 画图应用 \quad B. CLI 聊天机器人 \quad C. 视频播放器 \quad D. 电子表格

  \item 服务器先实现哪两个 tool？\\
  A. read / update document \quad B. search / delete \quad C. summarize / translate \quad D. login / logout

  \item mcp dev 的用途是什么？\\
  A. 运行前端 \quad B. 启动 Inspector 调试 server \quad C. 编译文档 \quad D. 安装依赖

  \item Inspector 能帮助你做什么？\\
  A. 只看日志 \quad B. 测试资源、提示词和工具 \quad C. 写前端样式 \quad D. 训练模型

  \item 项目中的文档最初存放在哪里？\\
  A. 远程数据库 \quad B. 内存 \quad C. 浏览器 localStorage \quad D. GitHub

\end{enumerate}

\subsection*{B. 判断题（3题）}
\begin{enumerate}
  \item 这个项目一开始就依赖完整的持久化存储。（\quad）
  \item 工具是 server 侧实现出来、供模型调用的能力。（\quad）
  \item Inspector 的目标之一是帮你验证工具返回值是否符合预期。（\quad）
\end{enumerate}

\subsection*{C. 简答题（2题）}
\begin{enumerate}
  \item 为什么课程会选择先做一个小型 CLI 项目？
  \item Inspector 能帮你发现哪些类别的问题？
\end{enumerate}

\section*{6. 参考答案}
\addcontentsline{toc}{section}{6. 参考答案}
\subsection*{选择题答案}
1.B \quad 2.A \quad 3.B \quad 4.B \quad 5.B

\subsection*{判断题答案}
1. 错 \quad 2. 对 \quad 3. 对

\subsection*{简答题要点}
\begin{enumerate}
  \item 因为小项目能把 MCP 的抽象概念压缩到一个可观察的工程场景里，便于理解 client 和 server 的配合。
  \item 它可以帮助检查工具定义、返回数据结构、连接状态以及 server 是否按预期响应。
\end{enumerate}

\end{document}